\section{模型假设与符号说明}
\subsection{假设及适用依据}
假设输入图与官方固定配置正确给出算子、张量、周期、依赖和硬件参数；算子不能被拆成更细的可调度原子，张量大小与算子周期在一次评价内确定。忽略 CPU 主控本身未由官方评价器计入的时间，不把网络传输、温度或动态频率引入模型。每个方案需满足题目规定的覆盖、无环、核内顺序、L1/UB 容量及资源占用约束。若容量压力触发溢出，新增 COPY 由原始核内启发式确定，不能由外层算法任意删除。

使用的 $100$ 张图由原附件逐字节无损提取，没有删除异常图、归一化周期或调整张量大小。导入时对全部图执行官方结构校验，并保存源文件与每份模型输入的 SHA-256；所有问题读取同一处理目录。由此，三问结果的差异来自场景语义及合法方案，而非改写数据。前期挑选小图和反例用于调试与同预算对照；最终汇总只以完整百图为准，不把开发样本当独立盲测集。

\subsection{固定硬件量}
L1 容量为 $524\,288$ bytes，UB 容量为 $131\,072$ bytes，共享 DDR 总带宽为 $60$ bytes/cycle。场景 A 的跨核任务等待为 $1000$ cycles，同核连续 Task 间隔为 $100$ cycles；场景 B 的跨核 COPY 释放延迟为 $500$ cycles。问题三的共享只读 L2 容量为 $1\,048\,576$ bytes，带宽为 $250$ bytes/cycle。以上均是主实验的固定配置，灵敏度扫描仅用于解释机制，不替换正式结果。
